\documentclass[12pt]{article}

\usepackage{amsmath,amssymb}
\usepackage[margin=1in]{geometry}
\usepackage{hyperref}
\usepackage{booktabs}
\usepackage{microtype}

% ── notation shortcuts ─────────────────────────────────────────────────────────
\newcommand{\chisq}{\chi^2}
\newcommand{\amp}{\mathcal{A}}
\newcommand{\iamp}{\mathcal{I}}

\title{FastSpecFit Bayesian Photometric SED Fitting}
\author{John Moustakas%
  \thanks{Grid design and fitting-algorithm derivations assisted by
    Claude Code (Claude Sonnet 5, Anthropic).}}
\date{\today}


\begin{document}
\maketitle

\section{Introduction}
\label{sec:intro}

This note documents \texttt{fastbayes}, \textsc{FastSpecFit}'s grid-based
Bayesian mode for fitting broadband photometry alone (no spectroscopy). It
is intended for samples where only photometry is available, or as a fast,
fully marginalized alternative to the gradient-based continuum fit in
\texttt{fastphot}. Rather than optimizing a handful of continuous stellar
population parameters, \texttt{fastbayes} evaluates every object against a
large, pre-synthesized grid of fully specified single stellar population
(SSP) models --- stellar continuum, nebular emission, dust attenuation, and
dust re-emission all baked in by FSPS --- and builds up the full discrete
posterior over the grid by likelihood weighting. The grid is built once by
\texttt{bin/build-bayesian-templates} (raw model spectra) and
\texttt{bin/build-bayesian-photometry} (synthesized broadband photometry on
a redshift grid, with IGM attenuation and cosmological dimming already
applied); \texttt{fastbayes} itself never touches FSPS or reads a raw
model spectrum for more than the one refined template needed for a given
object (Section~\ref{sec:refine}).

\section{Grid Design}
\label{sec:grid}

Each grid point is a single-age SSP (\texttt{sfh=0}) computed with FSPS's
Kriek \& Conroy (2013) dust attenuation curve (\texttt{dust\_type=4}, with
\texttt{dust1=0} and \texttt{dust2} $\equiv\tau$ as the free diffuse optical
depth) and reradiated dust emission (\texttt{add\_dust\_emission=True},
Draine \& Li 2007 templates). Nebular line and continuum emission are
included directly in the spectrum. The grid is a full factorial design over
the seven axes in Table~\ref{tab:grid}: every combination of the seven axis
values exists, so the flattened grid's $N$-dimensional structure --- and the
flat index of any neighboring grid point --- is recoverable directly from
the per-axis values via index arithmetic, with no separate shape bookkeeping
required.

\begin{table}[h]
\centering
\renewcommand{\arraystretch}{1.4}
\begin{tabular}{@{}llrl@{}}
\toprule
Axis & Symbol & $N$ & Range (spacing) \\
\midrule
Age & $\log_{10}(\mathrm{age/yr})$ & 20 & 7.000 to 10.137 (linear) \\
Metallicity & $\log_{10}(Z/Z_\odot)$ & 4 & $-0.8$ to $0.2$ (linear) \\
Dust optical depth & $\tau$ (\texttt{dust2}) & 12 & 0.0 to 4.0 (linear) \\
Dust index & $n_{\rm dust}$ (\texttt{dust\_index}) & 8 & $-1.5$ to $0.4$ (linear) \\
Dust-emission $U_{\rm min}$ & $\log_{10}(U_{\rm min})$ & 2 & 0.1 to 15.0 (log) \\
Dust-emission $\gamma$ & $\gamma$ & 2 & 0.0 to 0.15 (linear) \\
Dust-emission $Q_{\rm PAH}$ & $\log_{10}(Q_{\rm PAH})$ & 3 & 0.1 to 7.0 (log) \\
\bottomrule
\end{tabular}
\caption{The seven Bayesian grid axes (\texttt{GRID\_AXIS\_COLUMNS} in
  \texttt{fastspecfit.fastbayes}), in full-factorial nested-loop order.
  Total grid size $20\times4\times12\times8\times2\times2\times3 = 92{,}160$
  models. Axes built log-uniform (age, metallicity, $U_{\rm min}$,
  $Q_{\rm PAH}$) are refined and reported in $\log_{10}$ space so that
  their formal uncertainties (Section~\ref{sec:refine}) are symmetric;
  the remaining axes ($\tau$, $n_{\rm dust}$, $\gamma$) are refined and
  reported in linear space. $\tau$ and $n_{\rm dust}$ (well constrained by
  the UV/optical/NIR SED shape alone) gained resolution and had their upper
  bounds widened to match Leja et al. (2017)'s fiducial prior range, so the
  dustiest galaxies (edge-on spirals, LIRG-like systems) aren't truncated;
  $U_{\rm min}$ and $\gamma$ (essentially unconstrained without far-IR
  photometry) instead collapsed to a coarse low/high pair spanning their
  Herschel-informed fiducial range; $Q_{\rm PAH}$ (constrained by mid-IR
  photometry alone) kept its resolution, with only its upper bound widened.}
\label{tab:grid}
\end{table}

\section{Fitting Methodology}
\label{sec:fit}

For an object at redshift $z$, the grid's precomputed maggies
$m_b(t, z)$ --- per unit stellar mass formed, in band $b$, for template $t$
--- are linearly interpolated in $z$ to the object's redshift, and converted
to flux density $F_b(t)$ using the effective wavelength of band $b$. Because
the model is linear in stellar mass, the $\chisq$-minimizing, non-negative
mass amplitude $\amp(t)$ for every template is available in closed form:
\begin{equation}
  \boxed{
  \amp(t) \;=\; \max\!\left(0,\;
  \frac{\sum_b \iamp_b\, f_b\, F_b(t)}{\sum_b \iamp_b\, F_b(t)^2}
  \right),
  }
  \label{eq:amplitude}
\end{equation}
where $f_b$ is the observed flux density and $\iamp_b$ its inverse variance
in band $b$. The corresponding $\chisq$ for every template,
\begin{equation}
  \chisq(t) \;=\; \sum_b \iamp_b\,\bigl[f_b - \amp(t)\,F_b(t)\bigr]^2,
  \label{eq:chi2}
\end{equation}
is evaluated for the entire grid at once (a single vectorized pass over all
92{,}160 templates), giving the discrete maximum-likelihood template
$t_{\rm best} = \arg\min_t \chisq(t)$ and, via Gaussian likelihood weighting,
\begin{equation}
  w(t) \;\propto\; \exp\!\left[-\tfrac{1}{2}\bigl(\chisq(t) - \chisq_{\rm min}\bigr)\right],
  \qquad \sum_t w(t) = 1,
  \label{eq:weight}
\end{equation}
the full discrete posterior over every grid parameter (and over the derived
quantities $\log_{10} M_\star$ and $\mathrm{SFR}$, obtained by scaling each
template's tabulated per-unit-mass values by $\amp(t)$). Point estimates ---
mean, mode, and the 25th/50th/75th percentiles --- are reported for every
parameter directly from this weighted posterior.

\paragraph{Formal uncertainties.} A formal uncertainty $\sigma_x$ is also
reported for every parameter $x$ --- the seven grid axes, $\log_{10}
M_\star$, $\mathrm{SFR}$, and the rest-frame luminosities and $D_n(4000)$
index of Section~\ref{sec:derived} --- as the square root of the weighted
second moment of that parameter's marginalized discrete posterior,
\begin{equation}
  \sigma_x \;=\; \sqrt{\sum_t w(t)\,\bigl[x(t) - \bar{x}\bigr]^2},
  \qquad \bar{x} \;=\; \sum_t w(t)\, x(t).
  \label{eq:sigma}
\end{equation}
Because this is a genuinely marginal quantity, it grows to reflect real
posterior structure --- multiple modes, skewness, parameter degeneracies
(e.g., age--dust--metallicity) --- that the local $\Delta\chisq=1$ curvature
of Section~\ref{sec:refine} cannot see, being confined to a single grid
corner. The per-template value $x(t)$ entering Equation~\ref{eq:sigma}
depends on how the parameter scales with the mass amplitude $\amp(t)$: for
the seven grid axes and the model $D_n(4000)$ index (intensive quantities,
unaffected by $\amp(t)$), $x(t)$ is simply the tabulated per-template value;
for $\log_{10} M_\star$ and the rest-frame luminosities of
Section~\ref{sec:derived} (both logarithmic and linear in $\amp(t)$),
$x(t) = \log_{10}[\amp(t)] + \log_{10}[y_{\rm permass}(t)]$, using each
template's tabulated per-unit-mass value $y_{\rm permass}(t)$ --- one extra
per-template array lookup, not a per-object model rebuild; for
$\mathrm{SFR}$ (linear in $\amp(t)$, but reported and weighted in linear
space so that $\mathrm{SFR}=0$ remains exact, see below), $x(t) =
\amp(t)\, \mathrm{sfr}_{\rm permass}(t)$ directly. Absolute magnitudes and
K-corrections (Section~\ref{sec:derived}) are the one exception: their
reported uncertainty comes instead from the real observed photometry, not
this marginalized SED-fit posterior.

\paragraph{SFR proxy.} Each grid point is an instantaneous single-burst SSP,
for which FSPS's native ``current'' SFR is identically zero at any lookback
time (it is only meaningful for continuous/tabular star formation
histories). \texttt{fastbayes} instead derives, per template, the SFR
averaged over the most recent 100~Myr directly from the burst age: since a
delta-function burst's entire formed mass lies within that window whenever
$\mathrm{age} \le 100$~Myr and none of it otherwise, this per-unit-mass value
is exactly $10^{-8}~M_\odot\,\mathrm{yr}^{-1}$ for $\mathrm{age} \le
100$~Myr and $0$ otherwise -- computed from the grid's tabulated ages alone,
with no need to rerun FSPS or rebuild the grid. This mirrors the 100~Myr
averaging window used for the main (tabular-SFH) templates in
\texttt{fastspecfit.continuum}.

\section{Sub-Grid Refinement}
\label{sec:refine}

Because the grid has finite resolution, the discrete maximum-likelihood
point $t_{\rm best}$ is refined to sub-grid precision before its parameters,
model spectrum, and $\chisq$ are reported. The refinement exploits the
grid's full-factorial structure: for axis $a$ with best-fit index $i_0$
(not at a grid edge), the three already-computed $\chisq$ values at indices
$i_0-1, i_0, i_0+1$ --- holding every other axis fixed at its
$t_{\rm best}$ index --- define an exact quadratic interpolant (three points
determine a unique parabola, so no least-squares fit is needed; see the
closed-form coefficients in \texttt{fastbayes.\_parabola\_vertex3}).
Recasting that parabola as
\begin{equation}
  \chisq(x) \;=\; \chisq_0 + \left(\frac{x - x_0}{\delta x}\right)^2
  \label{eq:parabola}
\end{equation}
gives the refined coordinate $x_0$ along axis $a$, plus a $\Delta\chisq=1$
curvature width $\delta x$ ($\delta x = 1/\sqrt{a}$ for
$\chisq = a x^2 + bx + c$) used only as a validity check on the fit (a
concave-down or degenerate parabola, i.e.\ $a \le 0$, invalidates the
refinement); $\delta x$ is not itself reported as a formal uncertainty ---
that instead comes from the marginalized-posterior second moment
$\sigma_x$ of Equation~\ref{eq:sigma}, which is not confined to a single
grid corner and so captures real posterior structure that this local
curvature cannot. The refinement is skipped entirely (falling back to the
discrete grid value) whenever $i_0$ is at a grid edge, the fit is
degenerate or concave-down, or the vertex $x_0$ falls outside the
three-point bracket.

\paragraph{N-linear interpolation.} The same per-axis fractional offset
$f_a \in [0,1)$ toward the nearer neighbor (zero when unrefined) feeds an
N-linear interpolation over the $2^7=128$ grid corners bracketing the
refined point: corner $\kappa$, defined by one bit per axis, carries weight
\begin{equation}
  w_\kappa \;=\; \prod_{a=1}^{7} \bigl[\kappa_a\, f_a + (1-\kappa_a)(1-f_a)\bigr],
  \qquad \kappa_a \in \{0, 1\},
  \label{eq:corner_weight}
\end{equation}
computed for all 128 corners at once via array broadcasting over the bit
pattern, rather than a nested loop. Corners with $w_\kappa \le 0$ (any axis
exactly on a grid point) are dropped and the remaining weights renormalized.
The refined model photometry, mass amplitude, and $\chisq$
(equations~\ref{eq:amplitude}--\ref{eq:chi2}) are then recomputed from the
$w_\kappa$-weighted combination of the corner templates' photometry, so the
reported \texttt{CHI2} corresponds to the same refined model as the reported
parameters, not the unrefined discrete grid point. This same refined model
photometry, scaled by the refined mass amplitude, is written to the
\texttt{METADATA} output as \texttt{FLUX\_SYNTH\_PHOTMODEL\_*} --- the
grid-interpolated synthetic photometry consistent with what $\chisq$ was
actually computed from, as opposed to a fresh resynthesis of the resampled
rest-frame spectrum below. The corner templates'
raw rest-frame spectra are combined the same way (weighted by $w_\kappa$ and
scaled by the refined mass amplitude) to build the one refined rest-frame
spectrum used for the derived quantities in Section~\ref{sec:derived} and
written to the \texttt{MODELS} output extension.

\section{Derived Quantities}
\label{sec:derived}

IGM attenuation and cosmological dimming for the \emph{grid} photometry are
already baked into the photometry grid by \texttt{bin/build-bayesian-photometry};
\texttt{fastbayes} only needs the IGM model and cosmology directly to redshift
the one refined rest-frame spectrum per object, for:
\begin{itemize}
\item \textbf{K-corrections and absolute magnitudes}
  (\texttt{KCORR*}/\texttt{ABSMAG*}), from the same routine used by
  \texttt{fastspec}/\texttt{fastphot}, which also returns the purely
  synthetic, non-K-corrected model magnitude (\texttt{ABSMAG*\_SYNTH}) and
  propagates the real observed photometry's own inverse variance into
  \texttt{ABSMAG*\_IVAR} --- the one reported uncertainty in this section
  \emph{not} drawn from the SED-fit posterior of Section~\ref{sec:fit};
\item \textbf{Rest-frame luminosities} at a fixed set of reference
  wavelengths (1450, 1500, 1700, 2800, 3000, 5100~\AA, plus 3.4, 12, and
  22~$\mu$m --- roughly WISE $W1$/$W3$/$W4$ --- as monochromatic
  IR/SFR indicators for samples with adequate IR photometric coverage), by
  linear interpolation of the refined rest-frame spectrum;
\item \textbf{The model $D_n(4000)$ index}, from the refined rest-frame
  spectrum directly.
\end{itemize}
The rest-frame luminosities and $D_n(4000)$ index depend only on a
template's intrinsic rest-frame spectral shape, not on redshift or observed
photometry, so \texttt{bin/build-bayesian-templates} precomputes each
template's $D_n(4000)$ index and per-unit-mass monochromatic luminosities
once, at grid-build time, and stores them as additional \texttt{METADATA}
columns (mirroring \texttt{sfr}/\texttt{mstar}). Their formal uncertainties
(Equation~\ref{eq:sigma}) then require no per-object model rebuild --- just
one extra pair of vectorized array operations over the full, already-cached
posterior weights $w(t)$, unlike the point estimates above, which still
require the one refined rest-frame spectrum. No Monte Carlo/uncertainty
propagation analogous to \textsc{FastSpecFit}'s \texttt{nmonte} machinery
exists for this mode; every uncertainty reported here comes instead from
either the marginalized discrete posterior or the real observed photometry,
as described above.

\section*{Acknowledgments}

The grid-based fitting approach and sub-grid refinement scheme described
here were developed for \textsc{FastSpecFit}'s Bayesian photometric mode.

\end{document}
